\section*{Hoofdstuk \ref{ch:g8:combustion-and-fuels} --- Verbranding: brandstoffen, reactieproducten en broeikasgassen}

\begin{solution}{exo:g8:combustion-and-fuels:1}
\ce{C + O2 -> CO2}.
\end{solution}

\begin{solution}{exo:g8:combustion-and-fuels:2}
Met genoeg zuurstof wordt alle koolstof koolstofdioxide (volledig). Met
te weinig wordt een deel koolstofmono-oxide of roet (onvolledig).
\end{solution}

\begin{solution}{exo:g8:combustion-and-fuels:3}
Het heeft geen kleur, geen geur en geen smaak: niemand merkt het op, en
toch is het giftig.
\end{solution}

\begin{solution}{exo:g8:combustion-and-fuels:4}
Fossiel: steenkool, aardgas, benzine. Hernieuwbaar: hout, biogas,
ethanol uit suikerbieten.
\end{solution}

\begin{solution}{exo:g8:combustion-and-fuels:5}
\ce{C3H8 + 5O2 -> 3CO2 + 4H2O}.
\end{solution}

\begin{solution}{exo:g8:combustion-and-fuels:6}
Koolstof: 2 \ce{CO2}; waterstof: 3 \ce{H2O}; zuurstof rechts $4 + 3 =
7$, waarvan er 1 uit het ethanol komt, dus 3 \ce{O2}:
\ce{C2H6O + 3O2 -> 2CO2 + 3H2O}.
\end{solution}

\begin{solution}{exo:g8:combustion-and-fuels:7}
Roet, dus koolstof. De \omterm{def:g4:what-a-fire-needs:burning}{verbranding} is onvolledig: de vlam heeft te weinig
zuurstof.
\end{solution}

\begin{solution}{exo:g8:combustion-and-fuels:8}
Ongeveer 339 ppm in 1980 en ongeveer 370 ppm in 2000: een stijging van
ongeveer 30 ppm.
\end{solution}

\begin{solution}{exo:g8:combustion-and-fuels:9}
De koolstof van hout heeft de boom een paar jaar of decennia geleden uit
de \omterm{def:g4:air-a-mixture-of-gases:air}{lucht} gehaald: als je het verbrandt, geef je dat koolstofdioxide
terug. De koolstof van steenkool zat miljoenen jaren onder de grond
opgeslagen: als je het verbrandt, voeg je nieuw koolstofdioxide toe aan
de \omterm{def:g4:air-a-mixture-of-gases:air}{lucht}.
\end{solution}

\begin{solution}{exo:g8:combustion-and-fuels:10}
\ce{2C + O2 -> 2CO}.
\end{solution}

\begin{solution}{exo:g8:combustion-and-fuels:11}
$(427 - 316)/316 = 111/316 \approx 0.35$: ongeveer \qty{35}{\percent}.
\end{solution}

\begin{solution}{exo:g8:combustion-and-fuels:12}
Eerst is er genoeg zuurstof: \ce{CH4 + 2O2 -> CO2 + 2H2O}. Naarmate de
zuurstof in de afgesloten kamer opraakt, brandt de vlam met te weinig
\omterm{def:g4:air-a-mixture-of-gases:air}{lucht}: \ce{2CH4 + 3O2 -> 2CO + 4H2O}. Koolstofmono-oxide, reukloos en
giftig, hoopt zich op in de kamer.
\end{solution}

\begin{solution}{pb:g8:combustion-and-fuels:1}
\textbf{1.} \ce{C3H8 + 5O2 -> 3CO2 + 4H2O}.

\textbf{2.} Links: 3 C, 8 H, 10 O. Rechts: 3 C, $4 \times 2 = 8$ H,
$3 \times 2 + 4 = 10$ O. Ze klopt.

\textbf{3.} \ce{2C3H8 + 7O2 -> 6CO + 8H2O}.

\textbf{4.} De \omterm{def:g8:combustion-and-fuels:complete}{volledige verbranding}: 5 zuurstofmoleculen per
propaanmolecuul, tegen $7/2 = 3.5$ voor de onvolledige.

\textbf{5.} De \omterm{def:g4:air-a-mixture-of-gases:air}{lucht} in de gesloten tent heeft al snel te weinig
zuurstof; de vlam krijgt er niet genoeg meer van en brandt onvolledig,
geel.

\textbf{6.} Koolstofmono-oxide: geen kleur, geen geur (en geen smaak).

\textbf{7.} GHS02, de vlam: zeer licht ontvlambaar. GHS04, de gasfles:
een gas dat onder druk staat.

\textbf{8.} Bijvoorbeeld: gebruik hem alleen in de open \omterm{def:g4:air-a-mixture-of-gases:air}{lucht}, nooit in
een tent of een afgesloten ruimte; houd hem uit de buurt van alles wat
kan \omterm{def:g4:what-a-fire-needs:burning}{branden}; verwissel het busje ver van elke vlam.

\textbf{9.} $132 \div 44 = 3$ keer zwaarder.

\textbf{10.} Van de zuurstof uit de \omterm{def:g4:air-a-mixture-of-gases:air}{lucht}, die in het koolstofdioxide
(en het water) terechtkomt.

\textbf{11.} Een \omterm{def:g8:combustion-and-fuels:fossil}{fossiele brandstof} (het komt uit aardgas of aardolie):
het \omterm{def:g4:what-a-fire-needs:burning}{verbranden} ervan voegt koolstofdioxide toe aan de \omterm{def:g4:air-a-mixture-of-gases:air}{lucht}.

\textbf{12.} $450 \times 3 = \qty{1350}{g} = \qty{1.35}{kg}$
koolstofdioxide.
\end{solution}
